\chapter{مثال من الواقع: دورية تموينية على فرن وتنظيم ضبط}
\label{ch:real-example}

\section{القصة الحقيقية التي سنطبّقها}
\label{sec:real-story}
\textbf{لماذا؟} أتقنت في الفصل الأول بناء الكشف خطوة بخطوة. الآن نطبّق الأدوات نفسها على عمل تقوم به دوريات حماية المستهلك كل يوم: جولة على الأفران، ووزن ربطة الخبز أمام الزبون، وتنظيم ضبط عند المخالفة، ثم متابعة الضبط حتى التسوية أو الإحالة إلى القضاء.

في الواقع تخرج دورية من «دائرة حماية المستهلك» إلى الفرن. تزن ربطة الخبز، وتتحقق من الإعلان عن السعر، والنظافة والشروط الصحية، وتسجيل كميات الدقيق المستلمة على البطاقة الإلكترونية. إذا وجدت مخالفة تنظّم «ضبطاً تموينياً» وفق قانون حماية المستهلك الصادر بالمرسوم التشريعي رقم 8 لعام 2021. وإذا سُحبت عينة، تبقى المادة محجوزة حتى تظهر نتيجة التحليل. لبعض المخالفات يجيز القانون التسوية أمام المديرية خلال سبعة أيام من تنظيم الضبط، وإلا يُحال الضبط إلى القضاء.

\begin{table}[H]
\centering
\caption{العمل في الواقع وما يقابله في كيان}
\label{tab:real-map}
\begin{tabularx}{\linewidth}{R{3.2cm} Y Y}
\rcHeaderRow
\thh{المرحلة} & \thh{ما يحدث في الواقع} & \thh{الأداة في كيان}\\
تحديد الأفران & جولة على أفران حي أو منطقة، أو زيارة فرن بعينه بعد شكوى. & «خطة تكليف» على «المنشآت التجارية».\\
الزيارة & دورية من عضوين تعبئ محضراً واحداً. & نموذج «محضر دورية على فرن» بإجابة مشتركة.\\
الوزن والمخالفات & وزن الربطة، والسعر، والنظافة، والدقيق. & أسئلة اختيار، مع قواعد ظهور وإلزام.\\
الضبط والعينة & رقم ضبط، وعينة محجوزة عند الحاجة. & أسئلة تظهر وتُطلب فقط عند المخالفة.\\
تدقيق الضبط & رئيس الدائرة يراجع المحضر قبل متابعته. & خطوة موافقة بثلاث نتائج.\\
أثره على ملف الفرن & يُعرف الفرن بأنه مخالف بانتظار التسوية. & ربط ↑ كتابة إلى «الحالة الرقابية» عند الاعتماد.\\
التسوية أو الإحالة & موظف التسويات يتابع المهلة ويسجل النتيجة. & تعديل الحقل في الملف، ولوحة مؤشرات للمتابعة.\\
\end{tabularx}
\end{table}

\begin{rcnote}[ما الحقيقي وما التدريبي؟]
أسماء الوحدات، وعمل الدورية، ووزن الربطة أمام الزبون، والضبط، وحجز العينة، والتسوية خلال سبعة أيام لبعض المخالفات، والإحالة إلى القضاء: كلها مأخوذة من الواقع. أسماء الأشخاص وأرقام الضبوط والمواعيد تدريبية، وتحمل كلمة «تجريبي» أو «تدريب». لا نكتب وزناً معتمداً للربطة في الدليل؛ الدورية تأخذه من التعميم النافذ وقت الجولة.
\end{rcnote}
\textbf{تحقق:} تستطيع أن تفرق بين ثلاثة أشياء: المحضر (ما رأته الدورية)، والقرار (اعتماد رئيس الدائرة)، والمتابعة (التسوية أو الإحالة).

\section{أضف مناصب الدورية والتسويات}
\label{sec:real-positions}
\textbf{لماذا؟} الفصل الأول عمل بباحث ومدقق. الدورية الحقيقية لها رئيس ومراقب، ورئيس الدائرة يدقق الضبوط، وموظف مختص يتابع التسويات. نضيف هذه المناصب في «دائرة حماية المستهلك» نفسها، دون إنشاء وحدات جديدة.
\begin{enumerate}
\item افتح «الموارد البشرية» ← «وصف الوظيفة» ← «إضافة»، وأنشئ ثلاثة أوصاف: «رئيس دورية حماية المستهلك» \en{JD-CP-PATROL-LEAD}، و«مراقب تمويني» \en{JD-CP-MONITOR}، و«موظف التسويات والإحالة» \en{JD-CP-SETTLEMENT}.
\item افتح «المناصب» ← «إضافة»، وأنشئ المناصب من الجدول. الوحدة لكل منها «دائرة حماية المستهلك».
\item منصب «رئيس دائرة حماية المستهلك» موجود من الفصل الأول؛ لا تكرره. سنعيّن فيه شخصاً في الخطوة التالية.
\end{enumerate}
\begin{table}[H]
\centering
\caption{المناصب الجديدة}
\label{tab:real-positions}
\begin{tabularx}{\linewidth}{Y R{3.6cm} Y}
\rcHeaderRow
\thh{المنصب ورمزه} & \thh{الوصف الوظيفي} & \thh{المدير}\\
رئيس دورية حماية المستهلك — 01؛ \en{POS-CP-PATROL-LEAD-01} & \en{JD-CP-PATROL-LEAD} & رئيس دائرة حماية المستهلك.\\
مراقب تمويني — 01؛ \en{POS-CP-MONITOR-01} & \en{JD-CP-MONITOR} & رئيس دورية حماية المستهلك — 01.\\
موظف التسويات والإحالة — 01؛ \en{POS-CP-SETTLEMENT-01} & \en{JD-CP-SETTLEMENT} & رئيس دائرة حماية المستهلك.\\
\end{tabularx}
\end{table}
\textbf{تحقق:} افتح «مراقب تمويني — 01»؛ مديره رئيس الدورية، ووحدته دائرة حماية المستهلك.

\section{أضف الأعضاء وعيّن كل شخص في منصبه}
\label{sec:real-members}
\textbf{لماذا؟} التكليف والموافقة يصلان إلى شاغل المنصب. منصب بلا تعيين نشط لا يوصل العمل إلى أحد.
\begin{enumerate}
\item افتح «الأعضاء» ← «إضافة». أدخل الاسم و«رقم الموظف» وبريد عمل مختلفاً لكل عضو. اختر «مقعد كامل» لرئيسة الدائرة لأنها ستفتح لوحة المؤشرات، ويكفي «موظف» للبقية.
\item افتح «التعيينات» ← «إضافة». اختر العضو والمنصب، و«تاريخ البدء»، واترك «الحالة» «نشط»، ثم «إنشاء».
\item لا تغيّر تعيينات أحمد ونور وريم؛ يبقى مثال الفصل الأول يعمل كما هو.
\end{enumerate}
\begin{table}[H]
\centering
\caption{الأعضاء الجدد ومهامهم}
\label{tab:real-members}
\begin{tabularx}{\linewidth}{R{3.6cm} Y Y}
\rcHeaderRow
\thh{العضو ورقمه} & \thh{المنصب} & \thh{ما يفعله في المثال}\\
مازن قاسم — تجريبي؛ \en{DEMO-CP-004} & رئيس دورية حماية المستهلك — 01 & يقود الدورية ويرسل المحضر.\\
رنا شاهين — تجريبي؛ \en{DEMO-CP-005} & مراقب تمويني — 01 & تزن الربطة وتصوّر وتعبئ المحضر مع رئيس الدورية.\\
سلمى عبود — تجريبي؛ \en{DEMO-CP-006} & رئيس دائرة حماية المستهلك & تدقق الضبط وتعتمده أو تعيده أو ترفضه، وتتابع اللوحة.\\
فادي الخوري — تجريبي؛ \en{DEMO-CP-007} & موظف التسويات والإحالة — 01 & يسجّل نتيجة التسوية أو الإحالة في ملف الفرن.\\
\end{tabularx}
\end{table}
\textbf{تحقق:} في «التعيينات» أربعة تعيينات جديدة حالتها «نشط».

\section{وسّع ملف المنشأة بحقلين رقابيين}
\label{sec:real-fields}
\textbf{لماذا؟} نريد أن يعرف من يفتح ملف الفرن وضعه الرقابي الآن، وتاريخ آخر دورية عليه، دون فتح كل المحاضر. هذان الحقلان سيكتبهما المحضر بعد الاعتماد.
\begin{enumerate}
\item افتح «ادوات البناء» ← «الكيانات» ← «المنشآت التجارية». اضغط «إضافة حقل».
\item أضف حقلاً من نوع «اختيار» عنوانه «الحالة الرقابية»، بالخيارات الأربعة في الجدول. اترك «مصفوفة» متوقفاً؛ للفرن حالة واحدة في كل وقت.
\item أضف حقلاً من نوع «تاريخ» عنوانه «تاريخ آخر دورية».
\item اضغط «المراجعة والنشر» ثم «نشر … تغيير». الحقول لا تظهر للفريق قبل النشر.
\end{enumerate}
\begin{table}[H]
\centering
\caption{خيارات «الحالة الرقابية» ومتى تُستخدم}
\label{tab:real-status}
\begin{tabularx}{\linewidth}{R{4.4cm} Y}
\rcHeaderRow
\thh{الخيار} & \thh{من يضعه ومتى}\\
سليمة & المحضر المعتمد، عندما لا توجد مخالفة.\\
مخالفة — بانتظار التسوية & المحضر المعتمد، عندما يُنظَّم ضبط.\\
سُوّيت المخالفة & موظف التسويات، بعد التسوية ضمن المهلة.\\
أُحيل الضبط إلى القضاء & موظف التسويات، عند انتهاء المهلة دون تسوية أو عندما لا تجوز التسوية.\\
\end{tabularx}
\end{table}
\textbf{تحقق:} افتح «فرن الأمانة» من «حماية المستهلك» ← «المنشآت التجارية»؛ يظهر الحقلان فارغين. انسخ «القيمة المخزّنة» لكل خيار؛ ستحتاجها في الشروط.

\section{أنشئ نموذج «محضر دورية على فرن»}
\label{sec:real-form}
\textbf{لماذا؟} المحضر يحفظ ما رأته الدورية في هذه الزيارة. نبنيه بترتيب الزيارة الحقيقية: الفرن، ثم الوزن، ثم بقية البنود، ثم النتيجة والضبط، ثم الصور.
\begin{enumerate}
\item افتح «ادوات البناء» ← «النماذج» ← «إنشاء نموذج». اسم النموذج «محضر دورية على فرن»، و«كود النموذج» \en{002}، و«الاختصار المرجعي» \en{DWR}، والتطبيق «حماية المستهلك» والقسم «السجلات»، ثم «إنشاء».
\item أنشئ الأقسام والأسئلة من الجدول بالطريقة نفسها التي تعلمتها في الفصل الأول.
\item في سؤال «المخالفات المضبوطة» غيّر «يختار المجيب» ليقبل أكثر من خيار؛ قد يجمع الفرن عدة مخالفات في زيارة واحدة.
\end{enumerate}
\begin{table}[H]
\centering
\caption{أقسام المحضر وأسئلته}
\label{tab:real-questions}
\begin{tabularx}{\linewidth}{R{3cm} Y R{2.6cm}}
\rcHeaderRow
\thh{القسم} & \thh{السؤال ونوعه} & \thh{مطلوب؟}\\
بيانات الفرن & سؤال سجل «الفرن» من «المنشآت التجارية»، «موجود فقط». & دائماً\\
وزن الربطة & «الوزن المعتمد في التعميم النافذ (غرام)» رقم؛ «الوزن المقيس أمام الزبون (غرام)» رقم؛ «مطابقة الوزن» اختيار واحد: مطابق، ناقص. & دائماً\\
بنود الكشف & «الإعلان عن السعر» و«النظافة والشروط الصحية» و«تسجيل الدقيق على البطاقة الإلكترونية»، كل منها اختيار واحد: ملتزم، مخالف. & دائماً\\
النتيجة & «نتيجة الدورية» اختيار واحد: سليمة، مخالفة. & دائماً\\
الضبط & «المخالفات المضبوطة» اختيار متعدد؛ «رقم الضبط» نص؛ «هل سُحبت عينة؟» اختيار واحد: نعم، لا. & عند المخالفة\\
العينة & «رقم العينة» نص؛ «المخبر المرسلة إليه» نص. & عند سحب عينة\\
المرفقات & «صور الميزان والربطة» ملف، صور فقط. & دائماً\\
\end{tabularx}
\end{table}
\begin{rcnote}[لماذا نسجّل الوزن المعتمد في كل محضر؟]
الوزن المعتمد يحدده تعميم قد يتغير. إذا كتبناه قيمة ثابتة في النموذج صار كل محضر بعد التعديل خاطئاً. تسجيله في المحضر يحفظ المرجع الذي قيست عليه الربطة يوم الزيارة.
\end{rcnote}
\textbf{تحقق:} في «اختبار النموذج» تظهر الأقسام السبعة بالترتيب، ويمكن اختيار «فرن الأمانة».

\section{اجعل أسئلة الضبط تظهر عند المخالفة فقط}
\label{sec:real-rules}
\textbf{لماذا؟} الفرن السليم لا يحتاج رقم ضبط. إذا ظهرت أسئلة الضبط دائماً فقد يكتب المراقب أي شيء ليتجاوزها. القواعد تجعل المحضر يطلب ما يلزم فقط.
\begin{enumerate}
\item حدد سؤال «المخالفات المضبوطة»، ثم «مرئي» ← «يظهر عندما»: «نتيجة الدورية» «هو» «مخالفة». ثم «مطلوب» ← «مطلوب عندما» بالشرط نفسه. اضغط «حفظ في المسودة».
\item كرر الشرطين على «رقم الضبط» و«هل سُحبت عينة؟».
\item لسؤالي «رقم العينة» و«المخبر المرسلة إليه» اجعل الشرط «هل سُحبت عينة؟» «هو» «نعم».
\item أضف إلى «نتيجة الدورية» قاعدة تحقق بزر «+ قاعدة»: المستوى «تنبيه»، و«متى» «عند الإرسال»، والرسالة «اخترت سليمة مع بند مخالف. راجع البنود أو غيّر النتيجة». اجعل شرطها: «نتيجة الدورية» «هو» «سليمة» و«مطابقة الوزن» «هو» «ناقص». كرر الفكرة مع البنود الثلاثة إن أردت.
\end{enumerate}
\begin{rcwarn}[السؤال المخفي لا يُطلب]
لا تجعل «رقم الضبط» مطلوباً دائماً ثم تخفيه بشرط؛ النموذج يتجاوز السؤال المخفي ولا يطلبه. اجعل الظهور والإلزام بالشرط نفسه.
\end{rcwarn}
\textbf{تحقق:} في «اختبار النموذج» اختر «سليمة»؛ لا تظهر أسئلة الضبط. غيّر إلى «مخالفة» واترك «رقم الضبط» فارغاً وحاول الإرسال؛ يُمنع الإرسال. اختر «نعم» للعينة؛ يظهر سؤالا العينة مطلوبين.

\section{اربط المحضر بملف الفرن}
\label{sec:real-link}
\textbf{لماذا؟} بعد الاعتماد يجب أن يتغير ملف الفرن: حالته الرقابية وتاريخ آخر دورية. الربط بالكتابة يفعل ذلك دون أن يعدّل أحد الملف بيده.
\begin{enumerate}
\item افتح «السجلات» ← «من نوع سجل» ← «المنشآت التجارية»، واختر سؤال «الفرن» الموجود.
\item في «الحقول» حدد «الحالة الرقابية» و«تاريخ آخر دورية». اجعل «الربط» لكل منهما «↑ كتابة»، و«في النموذج» «مخفي»: الدورية لا تكتبهما بيدها.
\item حدد أيضاً «العنوان التفصيلي» بالربط «↓ قراءة» و«مقفل»، كي ترى الدورية عنوان الفرن للتأكد منه دون تعديله.
\item في «مراجعة» أضف الأسئلة إلى قسم «بيانات الفرن»، ثم «إضافة 3 أسئلة».
\item افتح سؤال «الحالة الرقابية» المخفي ← «القيمة الافتراضية» ← «إذا… وإلا»: إذا كانت «نتيجة الدورية» «سليمة» فالقيمة المخزّنة لخيار «سليمة»، وإلا القيمة المخزّنة لخيار «مخالفة — بانتظار التسوية». السلوك «يُحسب دائماً (مقفل)».
\item افتح سؤال «تاريخ آخر دورية» ← «القيمة الافتراضية» ← «مَن / متى» ← «اليوم / الآن».
\item من البطاقة اضغط «السماح في نوع السجل»، ثم «موافقة» على الحقلين. تأكد أن الكتابة «عند الاعتماد».
\end{enumerate}
\textbf{تحقق:} تظهر في بطاقة سؤال الفرن «يُكتب عند الاعتماد» و«مسموح · 2 حقول». المحضر لا يغيّر الملف قبل قرار رئيسة الدائرة.

\section{ضع المسار: رئيس الدائرة يدقق الضبط}
\label{sec:real-journey}
\textbf{لماذا؟} الضبط وثيقة قد تنتهي أمام القضاء. لذلك يراجعه رئيس الدائرة قبل أن يتغير ملف الفرن أو تبدأ متابعة التسوية.
\begin{enumerate}
\item افتح «المسار» ← «تفعيل المسار». غيّر «الاسم (للموظفين)» في «قيد المراجعة» إلى «قيد تدقيق الضبط».
\item اجعل سهم «إرسال» من «مسودة» ومن «بحاجة لاستكمال» إلى «قيد تدقيق الضبط».
\item أضف «خطوة موافقة» باسم «تدقيق محضر الدورية»، تعمل في «قيد تدقيق الضبط»، و«تبدأ» «عند الوصول إلى حالتها».
\item في «من يوقّع» ← «محدد العضوية» ← «المناصب» اختر «رئيس دائرة حماية المستهلك».
\item اضبط النتائج الثلاث كما في الجدول، واجعل «التعليق» «مطلوب» للرفض والإعادة.
\item اقرأ «الفحص» حتى تظهر «لا شيء يحتاج انتباهاً».
\end{enumerate}
\begin{table}[H]
\centering
\caption{قرارات رئيسة الدائرة}
\label{tab:real-decisions}
\begin{tabularx}{\linewidth}{R{3.2cm} R{3.2cm} Y}
\rcHeaderRow
\thh{النتيجة} & \thh{تنقل الطلب إلى} & \thh{متى تختارها}\\
اعتماد المحضر & مقبول & المحضر مكتمل وأدلته واضحة. يُكتب الحقلان في ملف الفرن.\\
إعادة للاستكمال & بحاجة لاستكمال & صورة الميزان غير واضحة، أو رقم الضبط ناقص.\\
رفض & مرفوض & المحضر نُظّم على منشأة خطأ أو بلا أساس. لا يتغير ملف الفرن.\\
\end{tabularx}
\end{table}
\textbf{تحقق:} على الخريطة سهم إرسال إلى «قيد تدقيق الضبط»، وثلاثة أسهم من الخطوة إلى «مقبول» و«مرفوض» و«بحاجة لاستكمال».

\begin{figure}[H]
\centering
\begin{tikzpicture}[node distance=0.5cm]
\node[stepBox] (a) {\RL{خطة التكليف: أفران الحي}};
\node[stepBox,below=of a] (b) {\RL{مازن ورنا: محضر مشترك}};
\node[stepBox,below=of b] (c) {\RL{سلمى: تدقيق الضبط}};
\node[stepBox,below=of c] (d) {\RL{ملف الفرن: الحالة الرقابية}};
\node[stepBox,below=of d] (e) {\RL{فادي: تسوية أو إحالة}};
\draw[flowsTo] (a) -- (b);
\draw[flowsTo] (b) -- (c);
\draw[flowsTo] (c) -- (d);
\draw[flowsTo] (d) -- (e);
\end{tikzpicture}
\rclegend
\caption{من الجولة إلى نتيجة الضبط}
\label{fig:real-route}
\shotnote{رسم توضيحي لتسلسل العمل، وليس لقطة من التطبيق. إذا أعادت سلمى المحضر للاستكمال يرجع إلى الدورية ثم يعود إليها. الفرن السليم يتوقف عند تحديث ملفه، ولا يصل إلى موظف التسويات.}
\end{figure}

\section{امنح الوصول وانشر المحضر}
\label{sec:real-access}
\textbf{لماذا؟} كل شخص يحتاج إذناً مختلفاً: الدورية تعبئ، ورئيسة الدائرة تقرأ وتقرر، وموظف التسويات يعدّل ملف الفرن بعد المهلة.
\begin{enumerate}
\item في النموذج افتح «الوصول» ← «إضافة أشخاص ثابتين»: مجموعة «الدورية» بمنصبي رئيس الدورية والمراقب، مع «قائمة» و«فتح» و«تعديل» و«إرسال الطلبات». ثم مجموعة «الدائرة» بمنصبي رئيس الدائرة وموظف التسويات، مع «قائمة» و«فتح».
\item افتح «الإعدادات» ← «الأدوار» ← «مراقب حماية المستهلك» (من الفصل الأول) ← «المسؤولين»، وأضف المناصب الأربعة الجديدة بنطاق «فقط». بهذا يقرؤون المنشآت.
\item أنشئ دوراً جديداً «متابعة تسويات الضبوط» وأسنده إلى «موظف التسويات والإحالة — 01». في «المنشآت التجارية» ← «الصلاحيات» ← «الأدوار» فعّل له «قائمة» و«فتح» و«تعديل»، ثم «المراجعة والنشر».
\item اختبر النموذج ثم اضغط «نشر» واكتب «الإصدار الأول لمحضر الدورية».
\end{enumerate}
\textbf{تحقق:} تظهر «تم نشر النموذج بنجاح». بحساب رنا يفتح النموذج ويظهر الفرن؛ وبحساب فادي يمكن تعديل «الحالة الرقابية» في ملف الفرن، ولا يمكن ذلك بحساب رنا.

\section{أطلق جولة على أفران الحي بإجابة مشتركة}
\label{sec:real-plan}
\textbf{لماذا؟} في الواقع لا تُكلَّف الدورية بفرن واحد، بل بجولة على أفران منطقة. وعضوا الدورية يكتبان محضراً واحداً، لا محضرين.
\begin{enumerate}
\item في النموذج افتح «إسناد»، وفعّل «يمكن التكليف به». «صلاحيات المكلَّف» «تعبئة وإرسال»، و«ينتهي التكليف عندما تصبح الإجابة» «مقبولة»، و«سؤال السجل» «الفرن»، و«الموعد» يوم واحد. اضغط «حفظ».
\item افتح «خطط التكليف» ← «خطة جديدة»، واختر «محضر دورية على فرن»، واسم الخطة «جولة أفران دمشق — تدريب»، ثم «إنشاء والإعداد».
\item «1 · السجلات»: «المنشآت التجارية»، و«كل السجلات المطابقة» بشرطين: «النشاط» «أفران ومخابز» و«المحافظة» «دمشق».
\item «2 · المرشحون»: «قائمة ثابتة» بمنصبي «رئيس دورية حماية المستهلك — 01» و«مراقب تمويني — 01».
\item «3 · التوزيع بين المرشحين»: «كلهم · إجابة مشتركة». هذا هو المحضر الواحد للدورية.
\item اضغط «تشغيل المعاينة» واقرأ عدد الأفران، ثم «حفظ» و«إنشاء التكليفات».
\end{enumerate}
\begin{rcnote}[شكوى على فرن واحد]
عندما تصل شكوى أو إخبار عن فرن بعينه، أنشئ خطة ثانية واختر في «السجلات» «اختيار سجلات» ثم الفرن المقصود فقط. المحضر نفسه والمسار نفسه.
\end{rcnote}
\textbf{تحقق:} الخطة «نشطة»، وتظهر الأفران المطابقة في «تكليفات مطلوبة منك» لدى مازن ورنا معاً، ولكل فرن إجابة واحدة.

\section{نفّذ الدورية ونظّم الضبط}
\label{sec:real-run}
\textbf{لماذا؟} هنا يلتقي الواقع بالنموذج: ما تقيسه الدورية أمام الزبون يصبح دليلاً محفوظاً.
\begin{enumerate}
\item بحساب رنا افتح تكليف «فرن الأمانة» واضغط «ابدأ». تأكد من العنوان المقفل أمامك.
\item اكتب الوزن المعتمد من التعميم النافذ، ثم الوزن الذي قسته أمام الزبون، واختر «مطابقة الوزن» «ناقص».
\item اختر «الإعلان عن السعر» «ملتزم»، و«النظافة والشروط الصحية» «ملتزم»، و«تسجيل الدقيق على البطاقة الإلكترونية» «مخالف».
\item اختر «نتيجة الدورية» «مخالفة». تظهر أسئلة الضبط: حدد المخالفتين، واكتب «رقم الضبط» \en{DWR-TRAIN-001}، واختر «هل سُحبت عينة؟» «لا».
\item ارفع صورة الميزان والربطة. بحساب مازن افتح المحضر نفسه وراجعه، ثم «إرسال الطلب».
\end{enumerate}
\textbf{تحقق:} المحضر «قيد تدقيق الضبط». ملف الفرن ما زال بلا حالة رقابية؛ الكتابة تنتظر الاعتماد.

\section{دقق الضبط واعتمده}
\label{sec:real-decide}
\textbf{لماذا؟} نجرّب الإعادة أولاً، لأن أكثر ما يُعاد في الواقع هو دليل ناقص.
\begin{enumerate}
\item بحساب سلمى افتح «موافقات بانتظارك» ثم المحضر. اقرأ «ما سيُنشأ أو يتغيّر في السجلات»: «سيتحدّث: فرن الأمانة».
\item اختر «إعادة للاستكمال» واكتب: «صورة الميزان لا تُظهر القراءة بوضوح. أرفقوا صورة أوضح».
\item بحساب رنا افتح المحضر من مهامها، وارفع صورة أوضح، ثم «إرسال الطلب».
\item بحساب سلمى اختر «اعتماد المحضر» وأكّد.
\end{enumerate}
\textbf{تحقق:} المحضر «مقبول». في ملف «فرن الأمانة»: «الحالة الرقابية» «مخالفة — بانتظار التسوية»، و«تاريخ آخر دورية» تاريخ اليوم. زر «السجلات · …» في المحضر يعرض «كُتبت».

\section{تابع المهلة وسجّل التسوية أو الإحالة}
\label{sec:real-settle}
\textbf{لماذا؟} اعتماد الضبط ليس نهاية العمل. لبعض المخالفات مهلة تسوية من سبعة أيام، وبعدها يُحال الضبط إلى القضاء. يجب أن يعرف الفريق كل يوم أي الأفران ما زالت معلقة.
\begin{enumerate}
\item افتح «ادوات البناء» ← «لوحات المعلومات» ← «إنشاء لوحة معلومات». الاسم «متابعة ضبوط الأفران»، والوصف «الأفران المخالفة بانتظار التسوية ونتائج الضبوط».
\item أضف مؤشر عدد على «المنشآت التجارية» مع مرشح داخل العنصر: «الحالة الرقابية» تساوي القيمة المخزّنة لـ«مخالفة — بانتظار التسوية». سمّه «بانتظار التسوية».
\item أضف رسماً يوزع الأفران حسب «الحالة الرقابية»، وجدولاً بأعمدة: اسم المنشأة، والمحافظة، والحالة الرقابية، وتاريخ آخر دورية.
\item انشر اللوحة وشاركها مع منصبي رئيس الدائرة وموظف التسويات.
\item بحساب فادي افتح ملف «فرن الأمانة». إذا سوّى الفرن المخالفة ضمن المهلة، غيّر «الحالة الرقابية» إلى «سُوّيت المخالفة». وإذا انتهت المهلة، أو كانت المخالفة لا تجوز تسويتها، اختر «أُحيل الضبط إلى القضاء». احفظ.
\end{enumerate}
\begin{rcnote}[التاريخ يحدد المهلة]
في الجدول رتّب الأفران حسب «تاريخ آخر دورية» من الأقدم. كل فرن مضى على ضبطه أكثر من سبعة أيام وهو «بانتظار التسوية» يحتاج قراراً اليوم.
\end{rcnote}
\textbf{تحقق:} ينقص مؤشر «بانتظار التسوية» بمقدار واحد بعد تعديل فادي، ويظهر الفرن في الرسم تحت حالته الجديدة.

\section{قائمة البطل: ما الذي يجب أن تستطيعه الآن}
\label{sec:real-hero}
\textbf{لماذا؟} تصبح مرجع فريقك عندما تجيب عن أسئلة العمل من النظام مباشرة، وتشرح سبب كل إعداد.
\begin{table}[H]
\centering
\caption{أثبت إتقانك}
\label{tab:real-hero}
\begin{tabularx}{\linewidth}{Y Y}
\rcHeaderRow
\thh{السؤال} & \thh{أين تجد الجواب؟}\\
كم فرناً مخالفاً ينتظر التسوية الآن؟ & مؤشر «بانتظار التسوية» في اللوحة.\\
ما الوزن الذي قيست عليه ربطة فرن الأمانة؟ & سؤال «الوزن المعتمد في التعميم النافذ» في المحضر نفسه.\\
لماذا أُعيد المحضر؟ & تعليق «إعادة للاستكمال» في سجل المحضر.\\
لماذا لم يطلب المحضر رقم ضبط لفرن سليم؟ & قاعدة الظهور والإلزام على أسئلة الضبط.\\
لماذا لم يتغير ملف الفرن قبل الاعتماد؟ & «يُكتب عند الاعتماد» في بطاقة سؤال الفرن.\\
لماذا تستطيع الدورية كتابة محضر واحد معاً؟ & «كلهم · إجابة مشتركة» في خطة التكليف.\\
\end{tabularx}
\end{table}
\begin{itemize}
\item كرر الجولة على فرن سليم، وتأكد أن الضبط لا يظهر وأن الحالة تصبح «سليمة».
\item جرّب محضراً بعينة: «هل سُحبت عينة؟» «نعم»، وتأكد من طلب رقم العينة والمخبر.
\item أنشئ خطة لشكوى على فرن واحد، واشرح لزميل الفرق بين الخطتين.
\end{itemize}
\textbf{تحقق:} أجبت عن الأسئلة الستة من النظام، ونفّذت تمريناً واحداً على الأقل دون الرجوع إلى الدليل.
